% The next line makes it so it compiles the main file if I hit the compile button
% when editing this file in TeXworks.
% !TEX root = ../AIDMA.tex
\chapter{Programming Fundamentals and Algorithms}\label{ch:algs}
The purpose of this chapter is to cover some of the 
basic programming concepts you 
may have picked up in previous classes while introducing you to some basic 
algorithms and new terminology that we will find useful as we continue our study of
discrete mathematics.
We will also practice our skills at writing proofs 
by sometimes proving that
an algorithm does as specified. 
 
Algorithms are presented in a syntax similar to Java and C++.  
This can be helpful since you may already be familiar with
one of these languages.  
On the other hand, this sort of syntax ties our hands more than one often likes when
discussing algorithms. When discussing algorithms, 
we often want to gloss over some of the implementation details.
For instance, we may not care about data types, or how parameters are passed
(i.e. by value or by reference), but by using a Java-like syntax we are forcing
ourselves to use particular data types and pass parameters in a certain way.  

Consider an algorithm that swaps two values (we will see
an implementation of this shortly).  The concept is the same regardless of what type of data is
being swapped.  But given our choice of syntax, we will give an implementation that assumes a
particular data type.  Most of the time the algorithms presented can be modified to work with
other data types.  

The issue of pass-by-value versus pass-by-reference is more complicated.  We will have a brief
discussion of this later, but the bottom line is that whenever you implement an algorithm from
{\em any} source, you need to consider how this and other language-specific features might
change how you understand the algorithm, how you implement it, and/or whether you even {\em can}.

\section{Basic Syntax and Algorithms}\label{section:algorithms}
There are various ways of defining what an algorithm is, and
there are subtle details that matter in certain contexts, but for
our purposes, the following simple definition will suffice.

\begin{df}
An {\bf algorithm}\index{algorithm} is a clear 
and unambiguous set of instructions
that accomplishes a task in a finite amount of time.
\end{df}

Before we see our first algorithm, we need some notation.

\begin{df}
\index{assignment operator}
The {\bf assignment operator}, $=$, assigns to the left-hand argument the value of the
right-hand argument.
\end{df}
\begin{exa}
The statement \verb@x = a + b@ means ``assign to $x$
the value of $a$ plus the value of $b$.''
\end{exa}

\begin{rem}
Most modern programming languages use $=$ for assignment.
Other symbols used include $:=$, $=:$, $<<$, $\leftarrow$, etc.

As it turns out, the most common symbol for assignment ($=$) is perhaps the most
confusing  for someone who is first learning to program. 
One of the most common assignment statements is \verb@x = x + 1;@.  
What this {\em means} is ``assign to the $x$ its current value plus one.''
However, what it {\em looks like} is the mathematical statement ``$x$ is equal to $x+1$'',
which is false for every value of $x$.  
If this has tripped you up in the recent past or still does,
fear not.  Eventually you will instinctively interpret it correctly, probably forgetting you
were ever confused by it.
\end{rem}
In our algorithms, we will work with various types of data,
but two will show up often and if you do not know a programming
language yet, they may be unfamiliar to you.
An \verb+int+ is a variable that stores an integer value 
(e.g. 45 or -123). A \verb+double+ or \verb+float+ is a
variable that stores a real number (e.g. 3.14159 or 
0.8675309). Technically, these store
{\em approximations} of real numbers, but we won't
concern ourselves with that detail here.
A \verb+boolean+ is a variable that can be either {\bf true} or {\bf false}.

In algorithms, {\em variables} store values that we can work with.

\begin{exa}
If I want to store a real number called x and have it store the value 5.5, I write it as
\begin{code}
double x = 5.5;
\end{code}
The ``\verb+double x+'' part of the code is specifying that I want \verb+x+ to store a real number. The ``\verb+x = 5.5+'' part is
specifying that I want to store the specific value 5.5 into \verb+x+.
\end{exa}

\begin{rem}
In many programming languages, semicolons (\verb+;+) 
are used at the end of each statement. Since this is not a
programming class, do not worry too much about this.
\end{rem}

Now we are ready for our first example of a simple algorithm. 
We will discuss the syntax after presenting the algorithm 
since it will make more sense if
we can use an example to discuss it.

\begin{exa}[Area of a Trapezoid]\label{exa:areaTrap}
Write an algorithm that gives the area of a trapezoid with height
$h$ and bases $a$ and $b$.  
\begin{sol}
One possible solution is
\begin{code}
double TrapezoidArea(double a, double b, double h) {
    return h*(a+b)/2;
}
\end{code}
\end{sol}
\end{exa}

The first line of most of our algorithms will give what is called
a {\em function prototype} that specifies 3 important pieces
of information about the algorithm:
\begin{enumerate}
\item What type of value/object is output by the algorithm.
We will often use the term {\em returned} instead
of output. They essentially mean the same thing.
\item The name of the algorithm.
\item The input to the algorithm, which we refer to as the 
{\em parameters}, given as a comma-separated list 
between parentheses after the name. 
Each parameter is given by specifying what type of
data it is and the name we will call it in the algorithm.
\end{enumerate}
In the previous example, the function prototype was
\begin{code}
    double TrapezoidArea(double a, double b, double h) {
\end{code}
We interpret it as follows:
\begin{enumerate}
\item
This algorithm returns a real number (since the return type
is \verb+double+). 
\item The name of the algorithm is \verb+TrapezoidArea+, which
makes sense because it is computing the area of a trapezoid.
\item
The inputs/parameters are 3 real numbers, \verb+a+,
\verb+b+, and \verb+h+.
\end{enumerate}
The curly brace (\verb+{+) at the end of the first line is used to denote that everything between that and the matching
end curly brace (\verb+}+) are the implementation of the algorithm.
Curly braces are used in general to denote
{\em blocks} of code. In addition to being used in this context,
they are also used for conditional statements and loops,
as we will see later. The simplest way to think about them is that
the are there to make it clear where a section (or block) of code
begins and ends.

The \verb+return+ keyword  is used
to indicate what value should be output/returned 
from the algorithm. For instance, if someone calls \verb+x=TrapezoidArea(a, b, h)+,
then $x$ will be assigned the value $h*(a+b)/2$ since this is what was {\em returned} by the algorithm.
Those who know {\em Java}, {\em C}, {\em C++}, or just about any other programming language 
should already be familiar with this concept.

\begin{exa}
If I execute the following code:
\begin{code}
double a = 10.0
double b = 5.0
double h = 3.5
double x = TrapezoidArea(a,b,h);
\end{code}
Then the variable \verb+x+ will have the value 
$(10.0+5.0)/2.5=4.2857142857.$
\end{exa}

\begin{exer}
Write an algorithm that returns the area of a square with sides of length $s$.
\begin{code}
double areaSquare(double s) {


}
\end{code}
\begin{answer}
\begin{code}
double areaSquare(double s) { return s*s; }
\end{code}
\end{answer}
\end{exer}


\begin{exa}[Swapping variables] \label{exa:swapping1}
\index{swap}
Write an algorithm that will interchange the values of two
variables, $x$ and $y$. 
That is, the contents of $x$ becomes that of $y$ and vice-versa.
\begin{sol}
We introduce a temporary variable $t$ in order to store
the contents of $x$ in $y$ without erasing the contents of $y$.
For simplicity, we will assume the data is of type \verb+Object+.
Here is the algorithm:
\begin{code}
  void swap(Object x, Object y) {
      Object t = x;  // Store x in a temporary variable
      x = y;      // x now has the original value of y
      y = t;      // y now has the original value of x
  }
\end{code}
\end{sol}
\end{exa}

\begin{rem}
Note that a return type of \verb+void+ means that
the algorithm does not return anything.
While some algorithms return values, other accomplish
some task and there is no need to return anything.
\end{rem}

It can be very useful to be able to prove that an algorithm actually does what we think it does.
Then when an error is found in a program we can focus our attention on the portions of the code
that we are uncertain about, ignoring the code that we {\em know} is correct.

\begin{exa}
Prove that the algorithm in Example~\ref{exa:swapping1} works correctly.
\begin{pf}
Assume the values $a$ and $b$ are passed into \verb+swap+.  Then at the beginning of the algorithm,
$x=a$ and $y=b$.  We need to prove that after the algorithm is finished, $x=b$ and $y=a$.

After the first line, $x$ and $y$ are unchanged and $t=a$ since it was
assigned the value stored in $x$, which is $a$.  After the second line, $x=b$ since it is assigned
the value stored in $y$, which is $b$.  Currently $x=b$, $y=b$, and $t=a$.  Finally, after the 
third line, $y=a$ since it is assigned the value stored in $t$, which is $a$.  Since $x$ is still $b$,
and $y=a$, the algorithm works correctly. 
\end{pf}
\end{exa}


\begin{rem}
The correctness of this algorithm (and a few others in this chapter)
is based on the assumption that the variables are 
{\bf passed by reference} rather than {\bf passed by value}.  

In C and C++, it is possible to pass by value or by reference, although we didn't use the
proper syntax to do so.  The \verb+*+ or \verb+&+ you sometimes
see in argument lists is related to this.
In Java, everything is passed by value and it is impossible to pass by reference.  
However, because non-primitive variables in Java are essentially reference variables 
(that is, they store a reference to an object, not the object itself), in some ways they act as if
they were passed by reference.
This is where things start to get complicated.
I don't want to get into the
subtleties here, especially since there are arguments about whether or not these are the best
terms to use.  Let me give an analogy instead.
\footnote{Inspired by a response on 
\url{http://stackoverflow.com/questions/373419/}.}

If I share a Google Doc with you, I am passing by reference.  
We both have a reference to the same document.
If you change the document, I will see the changes.  
If I change the document, you will see the changes.
On the other hand, if I e-mail you a Word document, I am passing by value.  
You have a copy of the document I have.  
Essentially, I copied the current {\bf value} (or contents) of the document
and gave that to you.
If you change the document, my document will remain unchanged.  
If I change my document, your document will remain unchanged.
This sounds pretty simple.  However, it gets more complicated.  
In Java, you can create a ``primitive'' Word document, 
but in a sense you can't create an ``object'' Word document.  
Instead, a Google Doc is created and you are given access (i.e. a reference) to it.
This is why in some ways primitive and object variables seem to act differently in Java.

I've already said too much.  
You will/did learn a lot more about this issue in another course.  
Here is the bottom line:  
The assumption being made in the various swap algorithms is that
when a variable is passed in, the algorithm has direct access to {\em that variable}
and not just a copy of it.  
Thus if changes are made to that variable in the algorithm, it is changing the variable that was
passed in.  This subtlety does not matter for most of the algorithms here.  
\end{rem}

\begin{rem}
We should be absolutely clear that it is {\bf impossible} 
to implement the \verb+swap+ method from Example~\ref{exa:swapping1} in Java.
In fact, there is no way to implement a method that swaps two arbitrary values in Java.
As we will see shortly, it {\bf is} possible to implement a method that swaps two elements from an array.
\end{rem}

\begin{rem}
One final note before we move on. Whether or not the \verb+swap+ method (or any method) 
can be implemented,
we can still use it in other algorithms as if it can.  This is because when discussing algorithms we are
usually more concerned about the {\bf idea} behind the algorithm, not all of the implementation details.
Using a method like \verb+swap+ in another algorithm often makes it easier to understand the overall
concept of that algorithm.  If we actually wanted to implement an algorithm that uses \verb+swap+ in a context where it is
impossible to implement, 
we would simply need to replace the call to \verb+swap+ with 
the 3 lines of code contained in it, replacing the variable names
as appropriate.
\end{rem}

\begin{ques}
Does the following swap algorithm work properly?  Why or why not?
\begin{code}
void swap(Object x, Object y) {
    x = y;
    y = x;
}
\end{code}
\answerlinetwo
\begin{answer}
It does not work.  To see why,
notice that if we pass in $a$ and $b$, then $x=a$ and $y=b$ at the beginning.
After the first line, $x=b$ and $y=b$. 
After the second line $x=b$ and $y=b$.  
The problem is that the first line overwrites the value stored in $x$ ($a$), and we can't recover it.
\end{answer}
\end{ques}

\begin{exa}\label{exa:swapnothird}\index{swap}
Write an algorithm that will interchange the values of two
variables $x$ and $y$ {\em without introducing a third variable}, assuming they are of some
numeric type.

\begin{sol}
The idea is to use sums and differences to store the
values. Assume that initially $x=a$ and $y = b$.
\begin{code}
  void swap(number x, number y) {
      x = x + y;  // x = a+b        and  y = b
      y = x - y;  // x = a+b        and  y = a+b-b = a  
      x = x - y;  // x = a+b-a = b  and  y = a
  }
\end{code}
Notice that the comments in the code sort of provide a proof that the algorithm is correct,
although keep reading for an important disclaimer.
\end{sol}
\end{exa}


\begin{exa}
It was mentioned that the comments in the algorithm from Example~\ref{exa:swapnothird}
provide a proof of its correctness.  What possibly faulty assumption is being made?
\begin{sol}
It is assumed that the arithmetic is performed with absolute precision, and that is 
not always the case.  
For instance, after the first line we are told that $x=a+b$.
What if $a=10,000,000,000$ and $b=.00000000001$?  Will $x$ really be {\em exactly} $10,000,000,000.00000000001$? If it isn't, the algorithm will not work properly. 

When talking about algorithms abstractly, we often ignore these
sorts of details, but is is extremely important to realize when we are
doing so.  
\end{sol}
\end{exa}

Problem~\ref{pro:twoVarSwap} explores whether or not the algorithm in Example~\ref{exa:swapnothird}
works for integer types--specifically 2's complement numbers.

Once again I want to emphasize that there can often be a huge 
difference between an algorithm as presented in a textbook and
the implementation of an algorithm in a programming language.
In this course, our focus is on the former. However, if you
become a programmer some day, it is extremely important
that you eventually understand some of the more subtle 
issues I have been pointing out.


\newpage

\section{Remainders and Rounding Revisited}\label{section:moddivRepeat}

Java, C, C++, and many other languages use $\%$ instead of the word {\em mod}.
For instance, you would write \verb+int x = a % n+ instead of \verb+int x = a mod n+.

\begin{rem}
Recall earlier that we mentioned that $a\bmod n$ has 
two possible values when $a<0$. 
When using the {\rm mod} operator in computer programs in situations where the
dividend might be negative, it is important to know which definition your
programming language/compiler uses.  {\em Java} returns a negative number when the dividend is negative. 
In {\em C}, the answer depends on the compiler.
\end{rem}

\begin{exer}
If you write a {\em C} program that computes $-45\bmod 4$, what are the two possible answers it might give you?\\
\answerline
\begin{answer}
Either $-1$ or $3$ are possible answers if we are uncertain whether it will return a positive
or negative answer.  But we know it is one of these.  It won't be -5, for instance.
\end{answer} 
\end{exer}

The next exercise explores a reasonable idea:  What if we want the answer to $ a\bmod b$ to 
always be between $0$ and $b-1$, even if $a$ is negative?  In other words, we want to
force the correct positive answer even if the compiler for the given language might return a
negative answer.

\begin{eval}\label{eval:negmod}
Although different programming languages and compilers might 
return different answers to the computation $a \bmod b$ when $a<0$, they always
return a value between $-(b-1)$ and $b-1$.  Given that fact, give an algorithm that
will always return an answer between $0$ and $b-1$, regardless of whether or not $a$ is negative.
Try to do it without using any conditional statements.

\begin{ssolenum}
\item
Use $(a\pmod{b} + b-1)/2$.  
Since it always returns a value between $-(b-1)$ and $b-1$ by adding $b-1$ to both sides
you get a value between $0$ and $2b-2$.  You then divide by $2$ to hit the target range of a return value that
is between $0$ and $b-1$ whether the number is positive or negative.\\
\evallinethree
\item
 Just return the absolute value of $a \bmod b$.\\
\evallinethree
\newpage
\item
Use the following:
\begin{lstlisting}
   int c = a % b;
   if(c<0) {
       return -c;
   } else {
       return c;
   }
   \end{lstlisting}
\evallinethree
\item
Use $(a\bmod b)\bmod b$.\\
\evallinethree
\end{ssolenum}
\begin{answer}
\begin{solevalenum}
\item
This solution is both incorrect and a bit confusing. 
The phrase `both sides' is confusing--both sides of what?
We don't have an equation in this problem.  
But there is a more serious problem. 
If you thought it was correct, go back and try to figure out why it is incorrect before you
continue reading this solution.
The main problem is that although this may return 
a value in the correct range, it doesn't always return the correct value.  In fact, what if
$(a\pmod{b} + b-1)$ is odd?  Mathematically, this would result in a non-integer result which is
clearly incorrect.  In most programming languages it would at least truncate and
return an integer--but again, not always the correct one.
This person focused on the wrong thing--getting the number in a particular range.  Although that
is important, they needed to think more about how to get the {\em correct} number.
They should have plugged in a few more values to double-check their logic.
\item
Incorrect. Generally speaking, $a\bmod b \neq -a\bmod b$.  
In other words, returning the absolute value when the result is negative is almost always incorrect.
\item
Incorrect.  If you think about if for a few minutes you should see that this is just 
one way of implementing the idea from the previous solution.
\item
Incorrect.  If $(a\bmod b)$ is negative, performing another $\bmod$ will still leave it negative.
\end{solevalenum}
\end{answer}
\end{eval}

\begin{exer}
Devise a correct solution to the Evaluate~\ref{eval:negmod}.\\
Answer:\\
\vspace*{1.75in}
\begin{answer}
There are several possible answers, but the slickest is probably:
\verb#(b+(a mod b)) mod b#.  Try it with both positive and negative
numbers for $a$ and convince yourself that it is correct.
\end{answer}
\end{exer}



\begin{eval}
Implement an algorithm that will round a real number \verb+x+ to the closest integer, 
but {\bf rounds {\em down} at .5}.
You can {\em only} use numbers, basic arithmetic (+, --, *, /), and \verb+floor(y)+ and/or 
\verb+ceiling(y)+ (which correspond to $\lfloor y \rfloor$ and $\lceil y\rceil$).
Don't worry about the data types (i.e. returning either a \verb+double+ or an \verb+int+ is fine
as long as the value stored represents an integer value).
\begin{ssolenum}
\item
\verb@ return floor(x+.49)@.\\
\evallinetwo
\item
\verb@ return floor(x+1/2)@.\\
\evallinetwo
\item
\verb@ return ceiling(x+.5)@.\\
\evallinetwo
\item
\verb@ return ceiling(x-.5)@.\\
\evallinetwo
\end{ssolenum}
\begin{answer}
\begin{solevalenum}
\item
This does not work.  What happens when $x=3.508$, for instance?
\item
This is incorrect for two reasons.  First, $1/2=0$ in most programming languages, so this
will {\em always} round down.
Second, even if we replaced this with $.5$ or $1.0/2.0$, it would round {\em up} at $.5$.
\item
Nice try, but still no good.  What if $x=2.5$?  This will round {\em up} to $3$.
Worse, what if $x=2.0001$?  Again, it rounds {\em up} to $3$ which is really bad.
\item
This one is correct.  Plug in values like $2$, $2.1$, and $2.5$ to see that it rounds down to
$2$ and values like $2.51$, $2.7$, and $2.9$ to see that it rounds up to $3$.
\end{solevalenum}
\end{answer}
\end{eval}




The floor function is important because in many programming languages, 
including Java, C, and C++, integer division {\em truncates}. That is, when you compute 
$n/k$ for integers $n$ and $k$, the result is rounded so it is {\em closer to zero} (as opposed to rounding down).
That means that if $n,k\geq 0$, $n/k$ rounds down to $\lfloor n/k\rfloor$.
But if $n<0$, $n/k$ rounds {\em up} to $\lceil n/k \rceil$.
So in Java, C, and C++, $3/4=-3/4=0$, $11/5=2$ and $-11/5=-2$, for instance.

\begin{exer}
Compute each of the following, assuming they are expressions 
in Java, C, or C++.

\vspace*{-.25in}
\begin{multicols}{3}
\begin{lenum}
\item 9/10 = \blankline{.5}
\item 10/10 = \blankline{.5}
\item 11/10 = \blankline{.5}
\item 14/10 = \blankline{.5}
\item 15/10 = \blankline{.5}
\item 19/10 = \blankline{.5}
\item 20/10 = \blankline{.5}
\item 90/10 = \blankline{.5}
\item -5/10 = \blankline{.5}
\item -10/10 = \blankline{.5}
\item -15/10 = \blankline{.5}
\item -20/10 = \blankline{.5}
\end{lenum}
\end{multicols}
\begin{answer}
(a) 0; (b) 1; (c) 1; (d) 1; (e) 1 (f) 1; (g) 2; (h) 9; (i) 0; (j) -1; (k) -1; (l) -2.
\end{answer}
\end{exer}

\newpage

\begin{eval}\label{exer:roundNotTruncate}
Let $n$ and $m$ be positive integers with $m>2$.  
Assuming integer division truncates, write an algorithm that will compute $n/m$,  
but will {\em round} the result instead of truncating it (round up at or above $.5$, down below $.5$).  
For instance, $5/4$ should return $1$, but 
$7/4$ should return $2$ instead of $1$.
You may only use basic integer arithmetic, not including the $\bmod$ operator.
\begin{ssolenum}
\item
\verb@floor(n/m+0.5)@\\
\evallinetwo
\item
\verb@floor((n/m) + 1/2)@\\
\evallinetwo
\item
\verb@(int) (n/m+0.5)@\\
\evallinetwo
\end{ssolenum}
\begin{answer}
\begin{solevalenum}
\item
$0.5$ is not an integer, and the \verb+floor+ function is not allowed.
\item
The \verb+floor+ function is not allowed.  Even if it were, this solution 
doesn't work. $1/2$ is evaluated to $0$ so it doesn't help.
\item
This one works, but $0.5$ is not allowed so it does not follow the directions.
\end{solevalenum}
\end{answer}
\end{eval}

Although the previous example may seem like it is based on an unnecessary restriction,
this is taken from a real-world situation.  
When writing code for an embedded device (e.g. a thermostat or microwave oven), code space
is often of great concern.  Performing just a single calculation using doubles/floats 
can add a lot of code since it needs to add certain code to deal with the data type.
Sometimes the amount of code added is too much since embedded processors have a lot
less space than the processor in your laptop or desktop computer.
Because of this,
some embedded programmers do everything they can to avoid all non-integer computations 
in their code when it is possible.  

\begin{exer}\label{exer:roundNotTruncatetwo}
Give a correct solution to round-instead-of-truncate problem from the previous example.\\
Answer:\\
\vspace*{.75in}
\begin{answer}
Two reasonable solutions include \verb@(n+m/2)/m@ and \verb@(2n+m)/(2m)@.
\end{answer}
\end{exer}

\newpage
\section{{\tt If-Else} Statements}\label{section:if}
{\em Conditional statements} allow an algorithm to 
do one of various things based on some condition.
In this section we introduce the most common 
{\em conditional statement}: the {\tt if-else} statement.

\begin{df} 
\index{if-else statement}
The {\bf if-else} control statement has the
following syntax:
\begin{code}
    if(expression) {
       blockA
    } else {
       blockB
    }
\end{code}

\noindent 
where \verb+expression+ must evaluate to a boolean value, 
and \verb+blockA+ and \verb+blockB+ are 0 or more statements
(so they can be blank).
If \verb+expression+ is true then the statements in \verb+blockA+ are executed. 
Otherwise, the statements in \verb+blockB+ are executed.
\end{df}

\begin{exa}\label{exa:maxTwo}\index{maximum!of two numbers}
Write an algorithm that will determine the maximum of two integers.  
Prove your algorithm is correct.
\begin{sol}
The following algorithm will work.
\begin{code}
  int max(int x, int y) {
      if(x >= y) {
          return x;
      } else {
          return y;
      }
  }
\end{code}
There are three possible cases.  If $x>y$, then $x$ is the maximum, and it is returned
since the algorithm returns $x$ if $x\geq y$.
If $x=y$, then they are both the maximum, so returning either is correct.  In this case
it returns $x$, the correct answer.
If $x<y$, then $y$ is the maximum and the algorithm returns $y$, which is the correct answer.
In any case it returns the correct answer.
\end{sol}
\end{exa}

\begin{exer}\index{maximum!of three numbers}
Write an algorithm that will determine the maximum of three numbers
that uses the algorithm from Example~\ref{exa:maxTwo}.
Prove that your algorithm is correct.
\begin{code}
  int max(int x, int y, int z) {
	
	
	
	
  }
\end{code}
\pflinethree
\begin{answer} 
Here is one possibility:
\begin{code}
    int max(int x, int y, int z) {
        int w = max(x,y);
        return max(w,z);
    }
\end{code}
We will use a proof by cases.  
\begin{itemize}
\item
If $x$ is the maximum, then $w=$\verb+max+$(x,y) = x$.
so it returns \verb+max+$(w,z)=x$, which is correct.  
\item
If $y$ is the maximum, the argument is essentially the same the previous case.
\item
If $z$ is the maximum, then $w$ is either $x$ or $y$, but in either case $w\leq z$, 
so it returns \verb+max+$(w,z)=z$.
\end{itemize}
In each case the algorithm returns the correct answer.
\end{answer}
\end{exer}

The previous exercise is an example of something that you are already familiar with: {\em code reuse}.
We {\em could have} written an algorithm from scratch, but sometimes it is easier to use one that already
exists.  Not only is the resulting algorithm simpler, it is easier to prove that it is correct
since we know that the algorithm it uses is correct.

\begin{exa}
Write an algorithm that determines whether a given integer is between two other integers. 
Have it return a \verb+boolean+ type 
(which is a variable type that can be either 
\verb+true+ or \verb+false+, as the name suggests).
\begin{sol}
Here is one possible solution:
\begin{code}
    boolean isBetween(int lo,int ho, in val) {
        if(lo < val && val < hi) {
            return true;
        else {
            return false;
        }
    }
\end{code}
Some people might prefer the conditional to be 
\verb+if(val > lo && val < hi)+, which is also perfectly fine.

In many programming languages, we can use
the following shorthand version.
\begin{code}
    boolean isBetween(int lo,int ho, in val) {
        return (lo < val && val < hi);
    }
\end{code}
This is because the result of the boolean expression 
turns out to be equal to what we want to return in each case.
That is, when \verb+lo < val && val < hi+ is \verb+true+,
we want to return \verb+true+, and if it is \verb+false+,
we want to return \verb+false+. So we can just ditch the
conditional statement in this case!
\end{sol}
\end{exa}

For simplicity, we will sometimes use \scode{print} to output results and not worry
about whitespace (i.e. spaces and newlines).  
Think of it as being equivalent to Java's \scode{System.out.print(i+" ")} or
C++'s \scode{cout<<i<<" "}, or C's \scode{printf("\%d ",i)} if you would like.

\begin{exa}
Here is an algorithm that prints the parity of an integer (that is,
whether it is even or odd):
\begin{code}
    void printParity(int n) {
        if(n%2==0) {
            print("Even");
        else {
            print("Odd");
        }
    }
\end{code}
\end{exa}

\newpage

\begin{exer} 
Write an algorithm that prints ``Hello'' if one enters a number between $4$ and $6$ (inclusive) and ``Goodbye'' otherwise. Use the function \verb+print(String s)+ to print.
You are not allowed to use any boolean operators like \verb+&&+, 
\verb+||+, etc.
\begin{code}
    void HelloGoodbye(int x) {
















    }
\end{code}
\begin{answer}
Here is a possible answer.
\begin{code}
    void HelloGoodbye(int x) {
        if(x >= 4) {
            if(x <= 6) {
                print("Hello");
            } else {
                print("Goodbye");
            }
        } else {
            print("Goodbye");
        }
    }
\end{code}
\end{answer}
\end{exer}


\begin{ques}
The solution given for the previous example uses three print statements, with two identical
print statements appearing in different places.  Is it possible to write the algorithm using
only two print statements while maintaining the restriction that you cannot use {\bf and} and {\bf or}?
If so, give that version of the algorithm.  If not, explain why not.\\
Answer:\\
\vspace*{2.5in}
\begin{answer}
It is possible.  If you thought it wasn't, go back and try to write the algorithm before
reading any further.

Here is one way to do it using an extra variable and an additional conditional statement.
\begin{code}
    void HelloGoodbye(int x) {
        boolean sayGoodbye = true;
        if(x >= 4) {
            if(x <= 6) {
                sayGoodbye = false;
            }
        }
        if(sayGoodbye) {
            print("Goodbye");
        } else {
            print("Hello");
        }
    }
\end{code}
It can also be done by using a return statement
(this version is not recommended!):
\begin{code}
    void HelloGoodbye(int x) {
        if(x >= 4) {
            if(x <= 6) {
                print("Hello");
                return;
            }
        }
        print("Goodbye");
    }
\end{code}
The second solution is simpler, 
but this sort of code (with somewhat random \verb+return+ statements in the middle
of them) can be tricky to debug if it is changed later.
Did you come up with a better solution than these?
\end{answer}
\end{ques}

\newpage

\begin{exer}
You need a program to execute some code only if the size of a list is not $0$.
The variable is named \verb+list+, and its size is \verb+list.size()+.
Give the expression that should go in the \verb+if+ statement.  
In fact, give two different expressions that will work.
\begin{enumerate}
\item \blanklinefill
\item \blanklinefill
\end{enumerate}
\begin{answer}
\verb+list.size()!=0+ and \verb+!(list.size()==0)+ are the most obvious solutions.
\end{answer}
\end{exer}

\begin{exa}
Write a code fragment that determines whether or not three numbers can
be the lengths of the sides of a triangle.
\begin{sol}
Let $a$, $b$, and $c$ be the numbers.  For simplicity, let's assume they are integers.
First we must have $a > 0$, $b>0$, and $c>0$. 
Also, the sum of any two of them must be larger than the third in order to form a triangle.
More specifically, we need $ a+b>c$, $ b + c>a$, and $c+a>b$.  
Since we need all of these to be true, this leads to the
following algorithm.
\begin{code}
IsItATriangle(int a, int b, int c) {
    if(a>0 && b>0 && c>0 && a+b>c && b+c>a && a+c>b) {
        return true;
    } else { return false; }
}\end{code}
\end{sol}
\end{exa}

\newpage
\section{Logic in programming}\label{section:logicInProgramming}

There are often times when conditional expressions can be simplified using logic. You are going to want to refer back
to Table~\ref{table:logicLaws} while reading this section.
Let's start with a simple example.

\begin{exa}
The conditional statement \verb+if(!(x<0))+ can be rewritten
as \verb+if(x>=0)+ to make it more readable.
\end{exa}

Too easy. Let's try a slightly more complicated one.

\begin{exa}
Simplify the following code.
\begin{code}
    if (x<=0 && (x>0 || x<=0)) {
       x=0;
    }
\end{code}
\begin{sol}
If we let $p$ be ``\verb+x<=0+,'' then $\neg p$ is ``\verb+x>0+,''
and the conditional statement is equivalent to  $p \wedge(\neg p\vee p)$.
But using negation and identity, we can see that
$p \wedge(\neg p\vee p)=p\wedge T=p$.
Thus, the code can be simplified to 
\begin{code}
    if(x<=0) { 
        x=0; 
    }
\end{code}
\end{sol}
\end{exa}

That wasn't so bad. Now we can do one that is a little trickier.
Pay attention on this one!
\begin{exa}
Simplify the following code as much as possible. 
\begin{code} 
    if ( !(x==0 || y<x) ) {
        x=y;
    }
\end{code}
\begin{sol}
First, notice that by DeMorgan's Law,  \verb+!(x==0 || y<x)+ is equivalent to
\verb+!(x==0) && !(y<x)+.  Simplifying a bit more, we get 
\verb+x!=0 && y>=x+.  Thus, the code becomes:
\begin{code}
    if (x!=0 && y>=x) {
        x=y
    }
\end{code}
This may not look much simpler, but it is definitely easier to understand.

This simplification can also be done by defining $p$ to be \verb+x==0+ and $q$ to be \verb+y<x+.
Then the expression is $\neg(p\vee q)$.  Applying De Morgan's Law, this is the same as $\neg p\wedge \neg q$,
which we translate back to \verb+!(x==0) && !(y<x)+ and simplify as in the final step above.
\end{sol}

\end{exa}
As the previous few examples demonstrate, you can apply the rules to propositions in various forms.
Sometimes it is useful to explicitly define $p$ and $q$ (and sometimes $r$) and write expressions using
formal mathematical notation, but at other times it is just as easy to apply the rules the the expressions
as they are.  In the previous example, we didn't gain that much by defining $p$ and $q$.  
But with more complicated expressions it certainly can be helpful.

\begin{rem}
A common mistake is to forget to use De Morgan's law when dealing with negation.
For instance, in the last example, replacing the code \verb+!(x==0 || y<x )+
with the code \verb+!(x==0) || !(y<x)+ would be incorrect.  
You cannot just distribute a negation among other terms.  
Always remember to use De Morgan's law: $\neg(p\vee q) \neq \neg p\vee\neg q$.
\end{rem}


\begin{exer}
Simplify the following code as much as possible.
\begin{code} 
    if ((x>0 && x<y) || (x>0 && x>=y)) {
        x=y;
    }
\end{code}
\vspace*{2.25in}

\begin{answer}
Let $p=$``$x>0$'' and $q=$``$x<y$''.  
Then the conditional above can be expressed as $(p\wedge q)\vee (p\wedge \neg q)$.
Using a few of the logical equivalences from Table~\ref{table:logicLaws}, 
it is not too difficult to see that
 $(p\wedge q) \vee (p\wedge \neg q) = p\wedge (q \vee \neg q) = p\wedge T = p.$ 
Therefore the code simplifies to:
\begin{code} 
    if (x>0) {
        x=y;
    }
\end{code}
\end{answer}
\end{exer}


\begin{eval}
Simplify the following code as much as possible. 
\begin{code} 
    if (x>0) {
        if(x<y || x>0) {
            x=y;
        }
    }
\end{code}
\begin{ssol}
Because the second if is in the first one which is if ($x>0$) then $x>0$ is duplicated but at
the same time to satisfy the second one we just need to keep the second if and cut the first one.
$x<y$ and $x>0$ are independent conditions so they cannot be more simplified.
So the answer is: 
\begin{code}
    if(x<y || x>0) {
        x=y;
    }
\end{code}
\end{ssol}
\evallinetwo
\begin{answer}
This is not equivalent to the original code.  Consider the case when $x=-1$ and $y=1$, for instance.
\end{answer}
\end{eval}

\begin{exer}
Simplify the following code as much as possible. 
\begin{code} 
    if (x>0) {
        if(x<y || x>0) {
            x=y;
        }
    }
\end{code}
\vspace*{1.5in}
\begin{answer}
After my head is done spinning trying to understand what they are
saying, I suspect something isn't right here. In fact,
since both conditions need to be true when if statements are nested, it is the same thing as a conjunction.
In other words, the two ifs are equivalent to \verb+if( x>0 && (x<y || x>0) )+.
By absorption, this is equivalent to \verb+if(x>0)+.
So the simplified code is:
\begin{code}
    if(x>0) {
        x=y;
    }
\end{code}
You can also think about it this way.  The assignment \verb+x=y+ cannot happen unless \verb+x>0+ due to the outer
if.\footnote{If you think about it, this is why the solution to this in the previous example failed.}
But the inner if has a disjunction, one part of which is \verb+x>0+, which we already know is true.
In other words, it doesn't matter whether or not \verb+x<y+.  This argument also leads to the solution we just gave.
\end{answer}
\end{exer}

Although some of these examples may seem a bit contrived, in some sense they {\em are} realistic.  As code is
refactored, code is added and removed in various places, conditionals are combined or separated, etc. and sometimes
it leads to conditionals that are more complicated than they need to be.  In addition, when working on large teams,
you will often work on code written by others.  Since some programmers don't have a good grasp on logic, you will
certainly run into conditional statements that are way more complicated and convoluted than necessary. 
As I believe these examples demonstrate, simplifying conditionals is not nearly as easy as one might think.
It takes great care to ensure that your simplified version is still correct.

\begin{rem}\index{short circuiting}
There is an important difference between the logical operators as discussed here and 
how they are implemented in programming languages such as Java, C, and C++.
It is something that is sometimes called {\bf short circuiting}.
You may be familiar with the concept even if you haven't heard it called that before.
It exploits the {\bf domination laws:}
$$F \wedge q = F$$
$$T \vee q = T$$
\end{rem}

\begin{exa}
Consider the following algorithm:
\begin{code}
if(x<0 || y<0) {
    x=y;
}
\end{code}
The conditional statement here involves an {\bf or} ($p\vee q$).
Notice that when $x<0$, then because of the {\bf or}, 
the conditional statement is true regardless of what the value of $y$ is. 
Most programming languages with therefore not even check the second part of the conditional in this case.

In general, when executing a conditional statement of the form
$p\vee q$,
$p$ is evaluated first. If it is true, the conditional is evaluated as true without determining the truth value of $q$. If $p$ is false,
then then $q$ is evaluated and its truth value used.
\end{exa}

\begin{exer}
Based on how short circuiting works with conditionals of the
form $p\vee q$, explain what it does with conditionals 
of the form $p\wedge q$.
\vspace*{1.75in}

\begin{answer}
$p$ is evaluated. If it is false, then the whole expression is evaluated 
as false without evaluating $q$ since no matter what the truth 
value of $q$ is, $p\wedge q=F$. On the other hand, if
$p$ is true, then the truth value of $p\wedge q$ is the same as
the truth value of $q$, so it evaluates $q$ and uses that as
the value of the whole expression.
\end{answer}
\end{exer}

\begin{rem}
If you forget that code short circuits, bad things can occasionally
happen.
Because of short circuiting, some of the code in your
conditional statement may never execute. 
The vast majority of the time that is fine.
However, there are rare cases when it might cause a problem.
Here is a simple example.
Say you want to increment both $x$ and $y$, and
then check if either is at least 100. 
Most sensible people would implement it like this:
{\rm
\begin{code}
    x++;
    y++;
    if( x==100 || y==100) {
        // do something
    }
\end{code}
}
That code will always work just fine. 
Alternatively, you can try to be clever and do this
(but don't! Keep reading):
{\rm
\begin{code}
    if( ++x==100 || ++y==100) {
        // do something
    }
\end{code}
}
First notice that the \verb@++@ comes before the variable,
so each variables is incremented before being compared to 100,
which is what we want. {\em However}, if $x$ becomes 100, then
short circuiting means that $y$ will never get incremented,
and that is incorrect!

Although this is a rather contrived example, there are certainly
cases where this comes up in real code, especially when dealing
with certain types of data structures.
\end{rem}

Below is a more complicated example involving short circuiting.
Earlier examples did not make demonstrate the
advantages to short circuiting, 
other than skipping the execution of a little code, but it is
very important in several contexts. The one in the next example 
involves checking if something exists before trying to access it.
This helps to avoid things like {\em segmentation faults} (C, C++), \verb+IndexOutOfBoundsException+ (Java),
and {\em null pointers} (many languages).

The example below uses a \verb+list+, which is exactly what it sounds like--an ordered list of 0 or more items of some sort 
(the list in the example contains integers).
If you have a list, calling \verb+list.isEmpty()+ will return
\verb+true+ if the list is empty and \verb+false+ if it is not empty.
Calling \verb+list.get(i)+ will get the $i$th element from the
list (where we use $0$ to get the first element, $1$ for 
the second, etc., for reasons I do not wish to get into right now). 
If there is no $i$th element, \verb+list.get(i)+ will crash
the program or return some bogus value, depending on 
how the list was implemented. This means it is important
to know that an item exists in a list before trying to access it.
This leads to the next example.

\begin{eval}
Rewrite the following segment of code so that it is as simple as possible and logically equivalent.
\begin{code}
    if( !(list.isEmpty() && list.get(0)>=100) && !(list.get(0)<100) ) {
        x++;
    } else {
        x--;
    }
    
\end{code}
\begin{ssolenum}
\item
The second and third statements mean the same thing.  Also, if the second is true then we got a value so
we know the list is not empty, so the first statement is unnecessary.  This leads to the following
equivalent code:
\begin{code}
    if(list.get(0) >= 100) {x++;} else {x--;}
\end{code}
\evallinetwo
\item
I used DeMorgan's law to obtain:
\begin{code}
    if(!list.isEmpty() || list.get(0) < 100) { 
        x++; 
    } else { 
        x--; 
    }
\end{code}
\evallinetwo
\item
Let $a$ be \verb+list.isEmpty()+ and $b$ be \verb+list.get(0)>=100+.  
But then $\neg b=$\verb+list.get(0)<100+.
The original expression is $\neg(a\wedge b)\wedge\neg(\neg b)$.
But 
\begin{eqnarray*}
\neg(a\wedge b)\wedge\neg(\neg b)&=&\neg(a\wedge b)\wedge  b
=(\neg a\vee \neg b)\wedge  b\\
&=&(\neg a\wedge  b)\vee(\neg b\wedge  b)
=(\neg a\wedge  b)\vee F
=\neg a\wedge  b
\end{eqnarray*}
So my simplified code is
\begin{code}
    if( !list.isEmpty() && list.get(0)>= 100 ) {
        x++;
    } else {
        x--;
    }
\end{code}
\evallinetwo
\end{ssolenum}
\begin{answer}
\begin{solevalenum}
\item
This solution is incorrect.
There are a few problems.  The obvious one is that the first statement actually prevents the program
from crashing so it is certainly not unnecessary!  Also, the second and third statements may be equivalent,
but how are they connected?  For instance, given the expression $\neg (A\wedge B)\vee \neg A$,
I cannot simply remove the `redundant' $\neg A$ to obtain an ``equivalent'' expression of $\neg (A\wedge B)$
(if necessary, plug in different truth values for $A$ and $B$ to convince yourself that these are not the same).
\item
This is not correct.  The second part of the expression seems to have disappeared.  
But how can we know it isn't equivalent?  We just need to find a scenario where the two versions do different
things. Notice that the `simplified' expression is true when the list is not empty regardless of the value of
element $0$.  But what if the list is not empty and element $0$ is 50?  The original expression is false and the
`simplified' expression is true.  Clearly not the same.
\item
This solution is not only correct, but it is very well argued.
\end{solevalenum}
\end{answer}
\end{eval}

\begin{ques}
In the solutions to the previous problem we said that the final solution was correct.
But there might be a catch.  Go back to the original code and the final solution
and look closer.  Is the final solution
{\em really} equivalent to the original?  Explain why or why not.\\
\evallinetwo
\begin{answer}
Technically speaking, the final solution is not equivalent.  
However, it turns out that it is {\em better} than the
original.  This is because the original code would actually crash if the list is empty.
Go back and look at the code and verify that this is the case.  
Then verify that the 
final simplified version will {\em not} crash.
\end{answer}
\end{ques}

Let's reinforce the point from the previous question.
The code from the third solution and the original code are
{\em logically} equivalent. 
However, they are not actually equivalent
when implemented in most programming languages. Why?
Because of short circuiting. 
What happens when you call \verb+get+ on an empty list? 
Does it return some bogus value? Or crash? Either way, 
the original code will either crash or return a bogus value.
However, the code from the third solution will always work 
correctly because short circuiting prevents the call to the
\verb+get+ that would cause the problem. 

\newpage

\section{Bitwise Operations}\label{section:bit}

In this section we will consider {\em bitwise operations}.  But first we need to review a few
concepts you are probably already familiar with.

Recall that a {\em boolean variable} is one that is either
true or false. 
Recall that a {\em bit} can have the value $0$ or $1$.  A bit can be used to represent a
Boolean variable by assigning $0$ to false and $1$ to true.
Table~\ref{tab:BooleanTruth} shows the truth tables for the various boolean operators 
that are available in many languages.
Notice that they are identical to the operators we discussed earlier except that we have 
replaced $T$ and $F$ with $1$ and $0$ and have used the notation from Java/C/C++ instead of the mathematical
notation.

\begin{table}[h]
\centering
$
\begin{array}{|cc||c|c|c|c|}
\hline  &  & AND & OR & XOR & IFF \\
\hline p & q & (p \&\& q) & (p || q) &p != q& (p  == q) \\
\hline 
1 & 1& 1 &1&0&1 \\
1 & 0& 0 &1&1&0 \\
0 & 1& 0 &1&1&0 \\
0 & 0& 0 &0&0&1 \\
\hline
\end{array}
$
\footnotesize\caption{Truth Tables for the Boolean Operators}\label{tab:BooleanTruth}
\end{table}

We don't usually think about \verb+!=+ being XOR 
and \verb+==+ being IFF (or biconditional).  We usually think of them in their more natural
interpretation: `not equal' and `equal'.  

\begin{rem}
A note of caution: Although Java is a lot like C and C++, how it deals with logical expressions
is very different.  Java has an explicit {\em boolean} type and you can only use the logical operators
on boolean values.  Further, conditional statements in Java require boolean values.   
In C and C++, the \verb+int+ type is used as a boolean value, where $0$ is false, and anything else is true.
This is very convenient, but can also cause some confusion.
\end{rem}

\begin{exa}
In C/C++, \verb+(5&&6)+, \verb+(5||0)+, \verb+(4!=5)+ are all true.  In Java the first two statements
are illegal. 
\end{exa}

Now it's time to extend the concept of boolean operators to integer data types (including
\verb+int+, \verb+short+, \verb+long+, \verb+byte+, etc.).

\begin{df}
A {\bf bitwise operation} is a boolean operation that operates on the individual bits of its argument(s).
\end{df}

\index{a bitwisenot@$\sim$ (bitwise compliment)}
\index{compliment, bitwise}
\index{bitwise operator!compliment}
\index{bitwise operator!NOT}
\begin{df}
The {\bf complement} or {\bf bitwise NOT}, usually denoted by \verb+~+, just flips each bit.
\end{df}

\begin{exa}
Assume \verb+10011001+ is in binary.  Then  
\verb+~10011001=01100110+. 
If this were a 32-bit integer, the answer would be
\verb+11111111111111111111111101100110+ since the leading 24 bits (which we assume to be $0$) 
would be flipped.  
\end{exa}

\begin{rem}
For simplicity, the rest of the examples will assume numbers are represented with 8 bits.
The concept is exactly the same regardless of how many bits are used for a particular data type.
\end{rem}

\begin{fib}
$255$ is \verb+11111111+ in binary. 
\verb+~11111111=00000000+, which is $0$ in decimal. 
Therefore, \verb+~255=0+.\\
Similarly, we can see that \verb+~240=15+ since $240$ is \blankline{1.5} in binary, and\\
\verb+~+\blankline{1.5} = \blankline{1.5}, which is \blankline{1} in decimal.
\begin{answer}
\verb+11110000+; \verb+11110000+; \verb+00001111+; $15$.
\end{answer}
\end{fib}

\begin{exer}
\verb+~11000110+=\blankline{2}
\begin{answer}
\verb+00111001+.
\end{answer}
\end{exer}


\index{a bitwiseand@\verb+&+ (bitwise AND)}
\index{AND (bitwise)}
\index{bitwise operator!AND}

\index{a bitwiseor@$\vert$ (bitwise OR)}
\index{OR (bitwise)}
\index{bitwise operator!OR}

\index{a bitwisexor@\verb+^+ (bitwise XOR)}
\index{XOR (bitwise)}
\index{bitwise operator!XOR}

\begin{df}
The following are the two-operator bitwise operators.
\begin{itemize}
\item
{\bf Bitwise AND},  denoted by \verb+&+, applies $\wedge$ to the corresponding bits
of each argument. 
\item
{\bf Bitwise OR},  denoted by \verb+|+, applies $\vee$ to the corresponding bits
of each argument. 
\item
{\bf Bitwise XOR},  denoted by \verb+^+, applies $\oplus$ to the corresponding bits
of each argument. 
\end{itemize}
\end{df}

We will present examples in table form rather than `code form' since it is much easier to see
what is going on when the bits are lined up.

\begin{exa}
\begin{tabular}[t]{rrrrr}
 			01011101&& 			  01011101			&& 01011101\\
\verb+&+ 	11010100	&&\verb+|+11010100	&&\verb+^+ 11010100\\ \cline{1-1}\cline{3-3}\cline{5-5}
			01010100&& 			  11011101			&& 10001001\\
 \end{tabular}
 \end{exa}
 

\begin{rem}
It is important to remember that \verb+&+ and \verb+&&+ are not the same thing!
The same holds for \verb+|+ and \verb+||+.
It is equally important to remember that \verb+^+ does {\em not} mean exponentiation in
most programming languages.
\end{rem}


\begin{exer}
{\Large
\begin{tabular}[t]{rrrrr}
 			11110000&& 			  11110000			&& 11110000\\
\verb+&+ 	11001100	&&\verb+|+11001100	&&\verb+^+ 11001100\\ \cline{1-1}\cline{3-3}\cline{5-5}
			&& 			  			&& \\
 \end{tabular}
 }
 \begin{answer}
 \verb+11000000+; \verb+11111100+; \verb+00111100+.
 \end{answer}
 \end{exer}
 
 \begin{rem}
Remember: boolean operators and the bitwise operators are different!
 \end{rem}


\newpage
\section{The {\tt for} loop}\label{section:for}

Here is the first of the two types of loops we will consider.

\begin{df}
\index{for loop}
\index{loop!for}
The {\bf for} loop has the following 
syntax:
\begin{code}
    for(initialize;condition;increment) {
       blockA
    }
\end{code}
where 
\begin{itemize}
\item
\verb+initialize+ is one or more statements that set up the initial conditions and is
executed once at the beginning of the loop.
\item
\verb+condition+ is the condition that is checked each time through the loop. It must evaluate to a boolean value.
If \verb+condition+ is true, the statements in \verb+blockA+ are executed followed by the code in \verb+increment+.
This process repeats until \verb+condition+ is false.
\item
\verb+increment+ is code that ensures the loop progresses.
It is executed at the end of every loop.
Typically \verb+increment+ is just a simple increment statement
(e.g. \verb@i++@), but it can be anything.
\end{itemize}
\end{df}

\begin{exa}\label{exa:factorial}\index{factorial}
Write an algorithm that returns $n!$  when given $n$.
\begin{sol}
Here is one possible algorithm.
\begin{code}
	int factorial(int n) {
	    if(n==0) {   return 1;
	    } else {
	        int fact = 1;
	        for(int i=1;i<=n;i++) {
	            fact = fact*i;
	        }
	        return fact;
	    }
	}
\end{code}
\end{sol}
\end{exa}

\begin{ques}
Does the \verb+factorial+ algorithm from Example~\ref{exa:factorial} ever do something unexpected?
If so, what does it do, when does it do it, and what should be done to fix it?\\
\answerlinethree
\begin{answer}
It will return 1 for negative numbers which does not really make sense.
%It will loop forever if $n<0$.  
This should be fixed in two ways.  
Since $n!$ is undefined when $n<0$, it can't return the correct answer for negative values.
So the first change is to have it return $-1$ if $n<0$.  
Some other value can be used, but $-1$ works well because $n!$ can't be negative, so if it returns $-1$, you know something is up.  
Second, this behavior should be clearly documented so that it clearly states that it returns $n!$ for $n\geq 0$, and $-1$ for $n<0$.
\end{answer}
\end{ques}

\newpage
\begin{eval}
Evaluate these algorithms that supposedly compute $n!$ for values of $n>0$. Don't worry about
what they do when $n\leq 0$.
\begin{ssolenum}
\item
\begin{code}

    int fact = 1;
    for(int i=0;i<=n;i++) {
        fact = fact*i;
    }
    return fact;
\end{code}
\evallinetwo
\item
\begin{code}

    int fact = 1;
    for(int i=2;i<=n;i++) {
        fact = fact*i;
    }
    return fact;
\end{code}
\evallinetwo
\item
\begin{code}

    int fact = 1;
    for(int i=n;i>0;i--) {
        fact = fact*i;
    }
    return fact;
\end{code}
\evallinetwo
\item
\begin{code}

    int fact = 1;
    for(int i=1;i<n;i++) {
        fact = fact*(n-i);
    }
    return fact;
\end{code}
\evallinetwo
\end{ssolenum}
\begin{answer}
\begin{solevalenum}
\item
This algorithm always returns $0$.  If you don't see it right away, carefully work through the algorithm
with a few values of $n$.
\item
This is correct.  It doesn't multiply by $1$, but that doesn't change the answer.
\item
This is also correct.  It is just multiplying the values in the reverse order of the other examples.
\item
This is incorrect.  It actually computes $(n-1)!$.
To fix this, does $i$ have to start at $0$, or go to $n$ (instead of $n-1$)?  
We'll leave it to you to work out which is correct.
\end{solevalenum}
\end{answer}
\end{eval}

\begin{exer}\label{exa:powers1} 
Write an algorithm that will compute $x^n$, 
where $x$ is a given real number and $n$ is a given positive integer.
\begin{code}
    double power(double x, int n) {
        
        
        
        
        
        
        
        
    }
\end{code}
\begin{answer}
This isn't really that much different than the algorithms to compute $n!$.
Here is one algorithm that does the job:
\begin{code}
    double power(double x, int n) {
        power = 1;
        for(int i=0;i<n;i++) {
            power = power*x;
        }
        return power;
    }
\end{code} 
The loop could also have been \verb@for(int i=1;i<=n;i++)@, 
or any loop that executes $n$ times since the loop index, $i$, is not used as part
of the calculation.
\end{answer}
\end{exer}


\newpage

\section{Arrays}\label{section:arrays}
\begin{df}
\index{array}
An {\bf array} is an aggregate of homogeneous types. The {\bf
length} of the array is the number of entries it has.
\end{df}
A $1$-dimensional array is akin to a mathematical vector. Thus
if $X$ is $1$-dimensional array of length $n$ then
$$X = (X[0], X[1], \ldots , X[n-1]).$$ 
You access the element at position $i$ of 
an array using the syntax $X[i]$.
We will
follow the convention of common languages like Java, C, and C++ by indexing the arrays starting from $0$.
This means that the last element is $X[n-1]$.
We will typically declare the length of the array at the beginning of
a code fragment by means of a comment.

\bigskip
A $2$-dimensional array is akin to a
mathematical matrix. Thus if $Y$ is a $2$-dimensional array with
$2$ rows and $3$ columns then 
$$Y = \begin{bmatrix} Y[0][0] &
Y[0][1] & Y[0][2] \cr Y[1][0] & Y[1][1] & Y[1][2] \cr
\end{bmatrix} .
$$
Here, you obtain the entry in row $i$ and column $j$ using
the syntax $Y[i][j]$.

In many programming languages
(including C++ and Java), arrays are declared using
the syntax 
\verb+int []a;+ (assuming an array of integers), 
where the \verb+[]+ indicates it is an array.
Without going into too much detail, there are actually two
parts to creating an array. The syntax \verb+int []a;+
is specifying that we want \verb+a+ to be an array, but
we need to allocate space for the array. To do that, we use
the syntax \verb+new int[n]+, where \verb+n+ is the 
desired size. Putting it all together, creating an array of
integers of size \verb+n+ is done using the syntax
\verb+int [] a = new int[n]+.
When specifying that we want to use an array as input
into an algorithm, we just use the \verb+int []a+ part since
we are not creating a new array.

\begin{exa}\index{maximum!array element}
Write an algorithm that determines the maximum element of a $1$-dimensional array of $n$ integers.
\begin{sol}
We declare the first value of the array (the $0$-th entry) to be the maximum (a {\em sentinel value}). 
Then we successively compare it to other $n - 1$ entries. 
If an entry is found to be larger than it, that entry is declared the maximum.
\begin{code}
MaxEntry(int[] X, int n) {
    int max = X[0];
    for(int i=1;i<n;i++) {
        if(X[i]>max) {
            max = X[i];
        }
    }
    return max;
}
\end{code}
\end{sol}
\end{exa}

If your primary language is Java, you might wonder why we did not use \verb+X.length+ in the previous
algorithm.  There are two reasons.  First, not all languages store the length of an array as part
of the array.  For example, C and C++ do not.  In these languages you always need to pass the length
along with an array.  Second, sometimes you want to be able to process only part of an array.  
Written as we did above, the algorithm will return the maximum of the first $n$ elements
of an array.  The algorithm works as long as the array has at least $n$ elements.


\begin{rem}
If an algorithm has an array and a variable $n$ as parameters, you can probably assume $n$ is
the length of the array unless it is otherwise specified.
\end{rem}

\begin{exer}
Give an algorithm that will return true if an array of integers either starts or ends with a $0$, and
false otherwise.  Assume array indexing starts at $0$ and that the array is of length $n$.
Use only one conditional statement.  Be sure to deal with the possibility of an empty array.
\begin{code}
    boolean startsOrEndsWithZero(int[] a, int n) {









    }
\end{code}
\begin{answer}
Here is one possible solution.  Note that the parentheses are necessary.
\begin{code}
    boolean startsOrEndsWithZero(int[] a, int n) {
        if( n>0 && (a[0]==0 || a[n-1]==0) ) {
            return true;
        else {
            return false;
        }
    }
\end{code}
\end{answer}
\end{exer}

\begin{ques}
Does the solution given for the previous exercise properly deal with arrays of size $0$ and $1$?  
Prove it.\\
\answerlinefour
\begin{answer}
The solution uses the expression  \verb+n>0 && (a[0]==0 || a[n-1]==0)+.
If $n=0$, the expression is false because of the \verb+&&+, so the algorithm returns false as it should
since an array with no elements certainly does not begin or end with a $0$.
If $n=1$, first note that $n-1=0$, so $a[0]$ and $a[n-1]$ refer to the same element.
Although this is redundant, it isn't a problem.  
If $a[0]=0$, the expression evaluates to $T\wedge (T\vee T)=T$, and the algorithm returns
true as expected.  
If $a[0]\neq 0$, the expression evaluates to $T\wedge (F\vee F)=F$, and the algorithm returns
false as expected.
\end{answer}
\end{ques}

\begin{exa}\label{exa:swapforArray}\index{swap}
Implement a method that swaps two elements of an array that works in Java and other languages
that can't pass by reference.
\begin{sol}
Here is a method that swaps two elements of an {\em integer} array.  Except for the
type of the parameter and \verb+temp+ variable, this works for any data type.
\begin{code}
swap(int[] X, int a, int b) {
    int temp = X[a];
    X[a]=X[b];
    X[b]=temp;
}
\end{code}
I don't want to get into the technical details of pass-by-value versus pass-by-reference since
that is really the subject of another course.  But briefly, 
this works because when the array is passed we have access to the individual array elements.
Therefore when we change them, they are changed in the original array.
\end{sol}
\end{exa}

\begin{exa}\index{reverse, an array}
An array $(X[0], \ldots X[n - 1])$ is given. 
Without introducing another array, put its entries in reverse order.
\label{exa:ReverseArray}
\begin{sol} 
Observe that we want to exchange the first and last element, the second and second-to-last element, etc.
That is, we want to exchange $X[0] \leftrightarrow X[n - 1]$, $X[1] \leftrightarrow X[n - 2]$, \ldots,
$X[k] \leftrightarrow X[n - k - 1]$.  But what value of $k$ is correct?  If we go all the way to $n-1$,
the result will be that every element is swapped and then swapped back, so we will accomplish nothing.
Hopefully you can see that if we swap elements when $k<n-k-1$, we will swap each element at most once.  
The ``at most once'' is because if the array has an odd number of elements, 
the middle element occurs when $k=n-k-1$, but we can skip it since it doesn't need to be swapped with anything.
Notice that $k<n-k-1$ if and only if $2k<n-1$. 
Since $k$ and $n$ are integers, this is equivalent to $2k\leq n-2$.
This is equivalent to $k\leq \lfloor (n-2)/2\rfloor$ by Corollary~\ref{floorCor}.
Thus, we need to swap the elements $0,1,\ldots, \lfloor (n-2)/2\rfloor$
with elements $n-1, n-2, \ldots, n-1-\lfloor (n-2)/2\rfloor=n-\lfloor n/2\rfloor$, respectively.
The following algorithm implements this idea.
\begin{code}
    reverseArray(int[] X, int n) {
        for(int i=0;i<=(n-2)/2;i++) {
            swap(X,i,n-i-1);
        }
    }
\end{code}
\end{sol}
\end{exa} 

\begin{ques}
The previous algorithm went until $i$ was $(n-2)/2$, not $\lfloor (n-2)/2\rfloor$.
Why is this O.K.?  Does it depend on the language?  Explain.\\
\answerlinetwo
\begin{answer}
As we already discussed, many languages truncate when performing integer division.
When the numbers are positive (as they are here), that is the same thing as taking the floor.
Even if a language does not do this, Theorem~\ref{floorThm} implies that it would still work.
\end{answer}
\end{ques}

\begin{ques}
Does the following algorithm correctly reverse the elements of an array? Explain.
\begin{code}
reverseArray(int[] X, int n) {
    for(int i=0;i<n/2;i++) {
        swap(X,i,n-i-1);
    }
}
\end{code}
\answerlinetwo
\begin{answer}
It is easiest to see that this is correct by comparing with the previous solution.
The only difference is that the condition went from $i<=(n-2)/2$ to $i<n/2$.
Notice that $(n-2)/2+1=n/2-2/2+1=n/2$.  
But since we also replaced $<=$ with $<$, it still stops at the same point.
\end{answer}
\end{ques}

Hopefully the previous example helps you realize that you need to be careful when working with arrays.
Formulas related to array indices change depending on whether arrays are indexed starting at $0$ or $1$.
In addition, formulas involving the number of elements in an array can be tricky, especially when the
formulas relate to partitioning the array into pieces (e.g. into two halves).
These can both lead to the so-called ``off by one'' error that is common in computer science.
The next example illustrates these problems, and one way to deal with it.


\begin{exa}
Give a formula for the index of the middle element of an array of size $n$.
If there are two middle elements (e.g. $n$ is even), use the first one.
\begin{sol}
Clearly the answer should be somewhere close to $n/2$.  Unfortunately, if $n$ is odd,
$n/2$ isn't an integer.  
And clearly the answer won't be the same when indexing starting at both $0$ and $1$.
 Maybe we should try a few concrete examples.\\

Let's first assume indexing starts at $1$.   
If $n=9$, the middle element is the 5th element, which has index $5 = \lceil 9/2\rceil$.  
If $n=10$, the middle element is also the 5th element.
Then the index is $5 = 10/2 =\lceil 10/2\rceil$.
Thus the formula $\lceil n/2\rceil$ should work.
You should plug in a few more values to convince yourself that this is correct.\\

Now let's assume indexing starts at $0$ (which is a lot more common).
There are a a few equivalent formulas we can come up with.
For starters, $\lceil n/2\rceil -1$ should work since this is just $1$ less 
than the answer above, and the indices are all shifted by one.
But let's come up with a formula from scratch.
If $n=9$, the index of the middle element is $4 = \lfloor 9/2\rfloor$.  
If $n=10$, the index is $4\neq \lfloor 10/2\rfloor$.  
So $\lfloor n/2\rfloor$ works when $n$ is odd, but not when $n$ is even.
This one is not as obvious as it was when we started indexing at $1$. 
With a little thought, you may realize that $\lfloor (n-1)/2\rfloor$ works.

\end{sol}
\end{exa}

\begin{ques}
The previous example seems to suggest that $\lceil n/2\rceil -1=\lfloor (n-1)/2\rfloor$
for all integers $n$.  Is this correct?  Do a few sample computations to try to convince
yourself of your answer.\\
\answerlinetwo
\begin{answer}
It is correct.  To convince yourself (but this is {\em not} a proof), 
plug the numbers $1$ through $5$ (or some other set of both even and odd values) into both sides
to see that you get the same number.
\end{answer}
\end{ques}

\begin{rem}
Always be very careful with formulas related to the index of an array.  
Double-check your logic by plugging in some values to be certain your formula is correct.
\end{rem}
%Before we jump into the next example, we  need introduce a few concepts.

\begin{exa}\label{exa:locker}[{\bf The Locker-Room Problem}] 

A locker room contains $n$ lockers, numbered $1$ through $n$.
Initially all doors are open. Person number $1$ enters and closes all the doors. Person number $2$ enters and opens all the doors
whose numbers are multiples of $2$.
 Person number $3$ enters and toggles all doors that are multiples of $3$.  That is, he closes them if they are
open and opens them if they are closed. 
This process continues, with person $i$ toggling each door that is a multiple of $i$.\\

Write an algorithm to determine
which lockers are closed when all $n$ people are done.
\begin{sol} \ \\
Here is one possible approach. We use a boolean array {\tt Locker} of size $n+1$ to denote the lockers (we will ignore \verb|Locker[0]|).  The value {\bf true} will denote an open
locker and the value {\bf false} will denote a closed
locker.\\
%\footnote{We will later see that those locker doors whose numbers are squares are the ones which are closed.}

\begin{code}
    LockerRoomProblem(boolean[] Locker, int n) {
        // Person 1: Close them all
        for(int i=1;i<=n;i++) {
           Locker[i]=false;
        }
        //People 2 through n: toggle appropriate ones
        for(int j=2;j<=n;j++) {
            for(k=j;k<=n;k++) {
                if(k%j==0) {
                    Locker[k] = !Locker[k];
                }
            }
        }
        // Print the results
        print("Closed:");
        for(int l=1;l<=n;l++) {
            if(Locker[l]==false) {
                print(l);
                print(" ");
            }
        }
    }
\end{code}
\end{sol}
\end{exa}
Can you see how to slightly improve the algorithm in Example~\ref{exa:locker}? 


Let's revisit short circuiting with a common use of it with arrays.
But first, in case you do not know, in languages like Java, you can get the length of an array \verb+a+ using \verb+a.length+
(unfortunately, this does {\em not} work in C or C++).
It can be helpful to determine the length before accessing 
an element of an array since \verb+a[i]+ will cause a
\verb+IndexOutOfBoundsException+ if \verb+i+$\geq$\verb+a.length+.

\begin{exa}\label{exa:ss}
Consider the statement \verb+if(x<a.length && a[x]!=0)+.  
The first domination law implies that when \verb+x>=a.length+, the expression in the if statement will evaluate to 
false regardless of the truth value of \verb+a[x]!=0+.  Therefore, many languages will simply
not evaluate the second part of the expression---they will {\em short circuit}, as we have discussed previously.
And it turns out that {\em this is a very good thing!} If
short circuiting did not occur, that code would crash 
whenever \verb+x>=a.length+ (If you do not see why,
look back and try to figure it out). 
That is one of the reasons many 
languages use short circuiting.

If \verb+x<a.length+, {\em only then} do we check that
\verb+a[x]!=0+, in which case everything is 
fine.\footnote{Well, sort of. If $x\leq 0$, this one will still crash.
You will fix this shortly!}
On the other hand, if \verb+x>=a.length+, the condition
is guaranteed to be false, so we short circuit and never check
the second condition {\em which is good, because it would
cause the code to crash!}

Therefore, even though the statements 
\verb+if(x<a.length && a[x]!=0)+ and 
\verb+if(a[x]!=0+ \verb+&&+ \verb+x<a.length)+ 
are {\em logically equivalent}, 
they are not {\em practically equivalent} (not an official term)
when implemented in a language that does short circuiting.

In general, since AND and OR operators are {\em commutative},
$p\vee q$ and $q\vee p$ are equivalent, as are
$p\wedge q$ and $q\wedge p$.
But the same statements implemented in a programming 
languages that short circuit are not precisely equivalent. 
But as this example shows, it is a very good thing since
(among other things) it allows us to write certain code 
more concisely without the possibility of crashing.

%% Too many examples that may or may not resonate,
%%especially non-cs students.
%The same thing happens for statements like \verb+if(x<0 || x>a.length)+. 
%When $x<0$, the expression is true regardless of the truth value of \verb+x>a.length+.  
%Again, many languages don't evaluate the second part of this expression if the first part is true.
%Of course, if the first part is false, the second part is evaluated since the truth value now depends
%on the truth value of the second part.
%
%There are at least two benefits of this.  First, it is more efficient since sometimes less code needs
%to be executed. 
%Second, it allows the checking of one condition before checking a second condition that might cause a crash.
%You may have used it in statements like the one above to make sure you don't index outside the bounds of
%an array.  Another use is to avoid attempting to access methods or fields when a variable refers to
%\verb+null+ (e.g. if \verb+a+ is a list, you might see code like \verb+if(a!=null && a.size()>0)+ which verifies that the list 
%actually exists before trying to access it).
%
%But this has at least two consequences that can cause subtle problems if you aren't careful.  
%First, although the AND and OR operators are {\em commutative} (e.g. $p\vee q$ and $q\vee p$ are equivalent),
%that is not always the case for boolean expressions in these languages.  
%For instance, the statement \verb+if(x<a.length && a[x]!=0)+ is not equivalent to \verb+if(a[x]!=0 && x<a.length)+ 
%since the first one always works, 
%but second one will cause a crash if $x<a.length$.
%Second, if the second part of the expression is code that 
%you expect will always be executed,
%you may spend a long time tracking down the bug that this creates.  
\end{exa}

\begin{ques}
Explain again: Why are \verb+if(x<a.length && a[x]!=0)+ and 
\verb+if(a[x]!=0+ \verb+&&+ \verb+x<a.length)+  
{\em logically equivalent}?
But then why are they not {\em practically equivalent}?\\
\answerlinethree
\begin{answer}
They are equivalent by the commutative property.
But they are not practically equivalent because
whenever \verb+x>=a.length+, the first statement works
properly because it does not try to access \verb+a[x]+
which avoids an \verb+IndexOutOfBoundsException+,
but the second one will try to access \verb+a[x]+ first,
causing an \verb+IndexOutOfBoundsException+ {\em before}
it checks if the index is valid.
\end{answer}
\end{ques}

\begin{exer}
Fix the conditional statement from Example~\ref{exa:ss}
so that is {\em never crashes}.

\vspace*{1in}
\begin{answer} This is simple---we just need to add a check
that the index is not negative:\\
\verb+if(x>=0 && x<a.length && a[x]!=0)+.
\end{answer}
\end{exer}

\begin{rem}
It turns out that if a language does not implement short circuiting,
we can obtain the same behavior by complicating the code just
a little bit---it involves using nested if-else statements instead of OR
and AND. You do not need worry too much about this since
most languages short circuit.
\end{rem}
%% Doesn't really fit here, and I rewrote the same example
%% (using DeMorgan's law) with simplified code earlier).
%{\bf Put the following example in appropriate contexts/locations.}
%In the following example, \verb+a+ is a list of some sort,
%\verb+a.size()+ returns how many things are on the list,
%and \verb+a.clear()+ removes everything from the list.
%Also, the code \verb+a==null+ is a way of asking whether or
%not \verb+a+ is
%
%\begin{exa}
%Simplify the following code as much as possible. 
%\begin{code} 
%    if ( !(a==null || a.size()<=0) ) {
%        a.clear();
%    }
%\end{code}
%\begin{sol}
%First, notice that by DeMorgan's Law,  \verb+!(a==null || a.size()<=0)+ is equivalent to
%\verb+!(a==null) && !(a.size()<=0)+.  Simplifying a bit more, we get 
%\verb+a!=null && a.size()>0+.  Thus, the code becomes:
%\begin{code}
%    if (a!=null && a.size()>0) {
%        a.clear();
%    }
%\end{code}
%This may not look much simpler, but it is much easier to understand.
%
%This simplification can also be done by defining $p$ to be \verb+a==null+ and $q$ to be \verb+a.size()<=0+.
%Then the expression is $\neg(p\vee q)$.  Applying De Morgan's Law, this is the same as $\neg p\wedge \neg q$,
%which we translate back to \verb+!(a==null) && !(a.size()<=0)+ and simplify as in the final step above.
%\end{sol}
%\end{exa}




%% No longer needed since logic comes earlier now.
%\begin{df} 
%\index{boolean!variable}
%A {\bf boolean variable} is a variable that can be either {\bf true} or {\bf false}.
%\end{df}
%\begin{df}
%\index{not operator}
%\index{negation operator}
%\index{operator!not}
%\index{operator!negation}
%\index{boolean!operator!not}
%\index{boolean!operator!negation}
%The {\bf not} unary operator changes the status of a boolean variable from {\bf true} to {\bf false} and vice-versa.  In Java, C, and C++, the {\em not} operator is $!$ and it appears before the expression being
%negated (e.g. \verb+!x+).
%\end{df}
%
%The {\em not} operator is essentially the same thing as the {\em negation} we discussed earlier.
%The difference is context---we are applying {\em not} to a boolean variable, whereas we applied 
%{\em negation} to a statement.

\newpage

\section{Quantifiers and Algorithms}\label{section:quantsAndAlgs}

Now let's see how quantifiers connect to algorithms.  
If you want to determine whether or not something (e.g. $P(x)$) is true for all values in a domain
(e.g., you want to determine the truth value of $\forall x P(x)$), 
one method is to simply loop through all of the values and test whether or not $P(x)$ is true.  
If it is false for any value, you know the answer is false.  
If you test them all and none of them were false, you know it is true.


\begin{exa}
Here is how you might determine if $\forall x P(x)$ is true or false for the 
domain $\{0, 1, 2, \ldots, 99\}$:
{\small
\begin{code}
    boolean isTrueForAll() {
        for(int i=0;i<100;i++) {
            if( !P(i) ) {
                return false;
            }
        }
        return true;
    }
\end{code}
}

Notice the negation in the code---this can trip you up if you aren't careful.
\end{exa}

\begin{exa}
Let $P(x)$ and $Q(x)$ be predicates and the domain be $\{0, 1, 2, \ldots, 99\}$.
What is \verb+isTrueForAll2()+ determining?  
\begin{code}
	    boolean isTrueForAll2() {
	        for(int i=0;i<100;i++) {
	            if( !P(i) && !Q(i) )
	                return false;
	        }
	        return true;
	    }
\end{code}
\begin{sol}
Notice that if both $P(i)$ and $Q(i)$ are false for the same value of $i$, 
it returns false, and otherwise it returns true.  Put another way, it returns
true if for every value of $i$, either $P(i)$ or $Q(i)$ is true. 
Thus, \verb+isTrueForAll2+ is determining the truth value of $\forall i(P(i)\vee Q(i))$.
\end{sol}
\end{exa}

\begin{exer}
Rewrite the expression \verb+( !P(i) && !Q(i) )+ from the previous example so that it uses only one negation.\\
Answer:\\
\vspace*{1in}
\begin{answer}
Using De Morgan's Law, we get \verb+!(P(i) || Q(i))+.
\end{answer}
\end{exer}

\newpage
\begin{exer}
Let $P(x)$ and $Q(x)$ be predicates and the domain be $\{0, 1, 2, \ldots, 99\}$.
What is \verb+isTrueForAll3()+ determining?
\begin{code}
		boolean isTrueForAll3() {
		    boolean result = true;
		    for(int i=0;i<100;i++) {
		        if(!P(i)) {
		            result = false;
		        }
		    }
		    if(result==true) {
		        return true;
		    }
		    for(int i=0;i<100;i++) {
		        if(!Q(i)) {
		            return false;
		        }
		    }
		    return true;
		}
\end{code}
\answerlinethree
\begin{answer}
First notice that if $P(i)$ is true for every value of $i$, \verb+result+ will be true at the
end of the first loop, so  \verb+isTrueForAll3+ will return true without even considering $Q$.
However, if $P(i)$ is false for any value of $i$, then it will go onto the second loop.
The second loop will return false if $Q(i)$ is false for any value of $i$.  But if $Q(i)$ is true
for all values of $i$, the method returns true.
So, how do we put this all together into a simple answer?
Notice that the only time it returns true is if either $P(i)$ is always true or if $Q(i)$ is always true.
In other words,  \verb+isTrueForAll3+ is determining the truth value of $(\forall i P(i)) \vee (\forall i Q(i))$.
By the way, notice that I used the variable \verb+i+ for both quantifiers. This is sort of like using
the same variables for two separate loops.
\end{answer}
\end{exer}

\begin{exa}
Now we are ready for the million dollar 
question:\footnote{There is no million dollars for answering this question.  It's just an expression.}
 Are \verb+isTrueForAll2+ and \verb+isTrueForAll3+ determining the same thing?
\begin{sol}
At first glance, it looks like they might be.  But we need to dig deeper, and we need to prove one way or the other.
To prove it, we would need to show that these expressions evaluate to the same truth value, regardless of what $P$ and $Q$ are.   
To disprove it, we just need to find a $P$ and a $Q$ for which these expressions have different truth values.
But let's first talk it through to see if we can figure out which answer seems to be correct.

$\forall i(P(i)\vee Q(i))$ is saying that for every value of $i$, either $P(i)$ or $Q(i)$ has to be true.
$\forall i P(i) \vee \forall i Q(i)$ is saying that either $P(i)$ has to be true for every $i$, or that
$Q(i)$ has to be true for every $i$.  These sound similar, but not exactly the same, so we cannot be sure yet.
In particular, we cannot jump to the conclusion that they are not equivalent because we described each with
different words.  There are many ways of wording the same concept. 

At this point we either need to try to tweak the wording so that we can see that they are really saying the same
thing, or we need to try to convince ourselves they aren't.  
Let's try the latter.

What if $P(i)$ is always true and $Q(i)$ is always false?  
Then  both statements are true.  But that doesn't prove that they are always both true, so this doesn't help.

Let's try something else.  
What if we can find a $P(i)$ and a $Q(i)$ such that for any given value of $i$, we
can ensure that either $P(i)$ or $Q(i)$ is true, but also that there is some value $j$ such that $P(j)$ is false
and some value $k\neq j$ such that $Q(k)$ is false?  
Then $\forall i(P(i)\vee Q(i))$ would be true, but $\forall i P(i) \vee \forall i Q(i)$ false, so this would work. 
But in order to be certain, 
we have to know that such a $P$ and $Q$ exist.\footnote{Consider this:  If I can find an even number
that is prime but is not $2$, then there would be at least 2 even primes.  That's great.  Unfortunately, 
I can't find such a number.}

What if we let $P(i)$ be ``$i$ is even'', $Q(i)$ be ``$i$ is odd'', and
the universe be $\BBZ$.  
Then $\forall i P(i) = \forall i Q(i) = F$, so $\forall i P(i) \vee \forall i Q(i)=F$, 
but  $\forall i(P(i)\vee Q(i))=T$. 
Now we have all of the pieces.  Let's put this all together in the form of a proof.

\begin{pf} (that  $\forall i(P(i)\vee Q(i)) \neq \forall i P(i) \vee \forall i Q(i)$)\\
Let $P(i)$ be ``$i$ is even'', $Q(i)$ be ``$i$ is odd'', and
the universe be $\BBZ$.  Then  $\forall i(P(i)\vee Q(i))$ is true since every integer is either even or odd.
On the other hand, $\forall i P(i)$ is false since there are integers that are not even 
and $\forall i Q(i)$ is false since there are integers that are not odd.  Thus, 
$\forall i P(i) \vee \forall i Q(i)$ is false.  Since they have different truth values, 
  $\forall i(P(i)\vee Q(i)) \neq \forall i P(i) \vee \forall i Q(i)$
\end{pf}
\end{sol}
\end{exa}

Let's do a slightly more practical example. 

\begin{exer}
Given an array $a$ of length $n$, give an algorithm that will determine the truth value of
$\exists x ( (a[x]\mbox{ is even }) \vee (a[x] \mbox{ is a multiple of 3}) )$.
\vspace*{2in}
\begin{answer}
\begin{code}
boolean has2MultipleOr3Multiple(int[] a, int n) {
   for(int i=0;i<n;i++) {
       if(a[i]%2==0 || a[i]%3==0) {
           return true;
       }
   }
    return false;
}
\end{code}
\end{answer}
\end{exer}

\newpage
\section{The {\tt while} loop}\label{section:while}
\begin{df}
\index{loop!while}
\index{while loop}
The {\bf while} loop has syntax:
\begin{code}
    while(condition) {
        blockA
    }
\end{code}
where \verb+condition+ evaluates to a boolean value and
\verb+blockA+ contains 1 or more statements.
If \verb+condition+ evaluates to true,
the statements in \verb+blockA+ will execute, 
and the condition will be checked again.
This repeats until \verb+condition+ evaluates to false.
\end{df}

While loops are generally used when you need to do something
some unknown number of times, usually because you need 
to check for some condition to be true before you know you are
done. The general rule of thumb for when to use \verb+for+ loops
versus \verb+while+ loops is that if you know the number of 
iterations, use a \verb+for+, and otherwise use \verb+while+.
On the other hand, anything that can be done with a \verb+for+
loop can be done with a \verb+while+ loop and vice-versa.
But using one of these loops in an unconventional manner can 
confuse people who read your code.

\begin{exa}
Here is a silly example of an algorithm to determine 
whether or not a number is a perfect square.
It is definitely {\em not} the slickest way to solve this problem.
\begin{code}
boolean isSquare(int x) {
   int i=1;
   while (i*i<=x) {
       if(i*i==x)  return true;
       i++
   }
   return false;
}
\end{code}
\end{exa}

The following is a common use of while loops: iterating over
a list of possibilities until you either find what you are looking
for (at which point you do something with it) or you don't
(which you know because you exited the while loop).

\begin{exa}
Here is a simple algorithm to determine if an array has any zero elements:
\begin{code}
boolean anyZeroes(int []a, int n) {
    int i=0;
    while(i<n) {
       if( a[i]==0 ) { 
           return true;
       }
       i++;
   }
   return false;
}
\end{code}
\end{exa}

\newpage

\begin{exer}
{\bf Sequential Search}\index{sequential search}
Implement the following algorithm that determines
whether or not a particular value is contained in the array.
It should return the index of the location of the value if present,
or -1 if not.
\begin{code}
int sequentialSearch(int []a, int n, int val) {









}
\end{code}
\begin{answer}
Here is one answer:\index{sequential search}
\begin{code}
int sequentialSearch(int []a, int n, int val) {
    int i=0;
    while(i<n) {
       if( a[i]==val ) { 
           return i;
       }
       i++;
   }
   return -1;
\end{code}
Notice that this algorithm is almost identical to the code from
the previous example---I just made 3 minor changes to the code!
This is because this is using a pattern that is common when 
dealing with arrays.
\end{answer}
\end{exer}


\begin{exa}
An array $X$ satisfies $X[0] \leq X[1] \leq \cdots \leq X[n - 1]$. 
Write an algorithm that finds the number of entries that are different.
\begin{sol}
Here is one possible approach.
\begin{code}
int differentElements(int[] X, int n) {
    int i = 0;
    int different = 1;
    while(i<n-1) {
        i++;
        if(x[i]!=x[i-1]) {
            different++;
        }
    }
    return different;
}
\end{code}
\end{sol}
\end{exa}

\begin{exer}
What will the following
algorithm return for $n = 5$? 
Trace the algorithm carefully, outlining all of your steps. 
\begin{code}
    int mystery(int n) {
        int x=0;
        int i=1;
        while(n>1) {
            if(n*i>4) {
                x=x+2*n;
            } else {
                x=x+n;
            }
            n=n-2;
            i++;
        }
        return x;
    }
\end{code}
%\answerlinetwo
\begin{answer} In the first iteration in the loop, 
$n = 5, i = 1$, $n*i > 4$ and thus $x = 10.$ 
Next, $n = 3$, $i = 2$, and we go through the loop again. 
Since $n*i > 4$, $x = 10 + 2*3 = 16$.
Finally, $n = 1, i = 3$, and the loop stops. Hence $x = 16$ is
returned.
\end{answer}
\end{exer}

An interesting example of the use of a while loop is 
implementing a simple algorithm to determine whether or
not a positive integer is prime.

\begin{thm}\label{thm:erathostenes_primality}
If $n\in\BBZ^+$, $n$ is either prime or has a prime factor  no greater than
$\sqrt{n}$.
\begin{pf}
If $n$ is prime there is nothing to prove. Assume that $n$ is composite. 
Then  $n$ can be written as the product $n = ab$ with $1 <a\leq b$, where $a$ and $b$ are integers. 
If every prime factor
of $n$ were greater than $\sqrt{n}$, then $a>\sqrt{n}$ and $b>\sqrt{n}$.
But then $n = ab > \sqrt{n}\sqrt{n} = n$, which is a contradiction. Thus $n$ must have a prime factor
no greater than $\sqrt{n}$.
\end{pf}
\end{thm}


\begin{exa}
To determine whether $103$ is prime we proceed as follows. 
Observe that $\lfloor \sqrt{103} \rfloor = 10$ 
(According to Theorem~\ref{floorThm}, we only need concern ourselves with the floor).
We now divide $103$ by every prime no greater than $10$. If one of these primes divides $103$, then $103$
is not a prime. Otherwise, $103$ is a prime. 
Notice that 
$103 \bmod 2 = 1$, $103 \bmod 3 = 1$, $103 \bmod 5 = 3$, and $103 \bmod 7 = 5$.
Since none of these remainders is $0$, $103$ is prime.
\end{exa}

\begin{exer}
Give a complete proof of whether or not $101$ is prime.\\
\pflinetwo
\begin{answer}
$\lfloor \sqrt{101}\rfloor=10$.  The only primes less than $10$ are $2$, $3$, $5$, and $7$.
Since $101\bmod 2=1$, $101\bmod 3 = 2$, $101\bmod 5 = 1$, and $101\bmod 7 = 3$, none of
which are $0$, Theorem~\ref{thm:erathostenes_primality} tells us that $101$ is prime.
\end{answer}
\end{exer}

\begin{exer}
Give a complete proof of whether or not $323$ is prime.\\
\pflinetwo
\begin{answer}
$323=17*19$, so it is not prime.  I determined this by seeing if any of the primes
no greater than $\lfloor\sqrt{323}\rfloor=17$ were factors.
Although $2$, $3$, $5$, $7$, $11$, and $13$, are not, $17$ is.
\end{answer}
\end{exer}


\begin{exa}\label{exa:primality}\index{primality testing}
How can we determine if an integer $n$ is prime?
If $n\leq 1$, it is not prime.
Otherwise, for each integer from $2$ to $\sqrt{n}$,
if $n$ is a multiple of it, $n$ is not prime.
If $n$ was not a multiple of any of these, it must be prime
by Theorem~\ref{thm:erathostenes_primality}.  
Here is the algorithm based on these ideas.
\label{exa:isprime1}
\begin{code}
boolean isPrime(int n) {
    if(n<=1) {  return false; }
    else {
        int i = 2;
        while( i <= sqrt(n) ) {
            if( n%i==0 ) { return false; }
            i++;
        }
        return true;  // It had no factors.
}   }
\end{code}
\end{exa}

\begin{exer}
Improve the algorithm from the previous example
by making it about twice as fast. Hint: The fact that
I am asking you to make it {\em twice} as fast is a hint.

\vspace*{3.5in}
\begin{answer}
No need to test the even integers. So
we can modify the code to skip checking even numbers,
so long as we first make sure that it gets the correct
answer for 2. If you didn't think about this,
go back and try to write the algorithm before
continuing to read the solution. 
Here is the modified code:
\begin{code}
boolean isPrime(int n) {
    if(n<=1) {      // Anything less than 2 is not prime
        return false; 
    } else if (n==2) {
       return true;    // Need to deal with this as a special case
    } else if( n%2==0 ) {  // Discard even numbers.
        return false;
    } else {
        // Determine if it has any odd factors.
        int i = 3;
        while( i <= sqrt(n) ) {
            if( n%i==0 ) {
                return false; 
            }
            i = i + 2;
        }
        return true;  // It had no factors.
    }
}
\end{code}
\end{answer}
\end{exer}

\begin{rem}
It should be noted that although this algorithms in the last
example and exercise work, 
they are not very practical for large values of $n$.
For instance, if we want to factor the number 
$12,867,530,934,567,891$,
we would still need to check about 
$\sqrt{12,867,530,934,567,891}\approx 113,435,140$ factors.

In fact, there is no known algorithm that can factor numbers efficiently on a ``classical'' computer.
The most commonly used public-key cryptosystems rely on the assumption that there is no efficient algorithm
to factor a number.  
However, if you have a quantum computer, you are in luck.  
Shor's algorithm actually {\bf can} factor numbers efficiently.
\end{rem}

Can we do the same thing for multiples of $3$, $5$, etc.
to make the algorithm even faster? 
That is, once we know that a number is not divisible by $3$, we don't really need
to ask if it is divisible by $6$, $9$, $12$, etc.  Similarly for 5, 7, 11, etc. Can we somehow make it skip all of these multiples as we go?
A first step would be to try to skip just multiples of either 2 or 3 (or both). If you can do that, you can then take into account 5, etc.
But does it generalize? And would it help? I will leave these as
questions for you to ponder.




\newpage

\begin{exer}
Use the fact that integer division truncates to write 
an algorithm that reverses the digits of a given positive integer. 
For example, if $123476$ is the input, the output should
be $674321$. 
You should be able to do it with one extra variable,
one {\tt while} loop, one $\bmod$ operation,
one multiplication by $10$, one division by $10$, and one addition.
\begin{code}
    int reverseDigits(int n) {











    }
\end{code}
\begin{answer} 
The following algorithm does the job.
\begin{code}
	int reverseDigits(int n) {
		x=0;
		while(n!=0) {
		    x = x*10+n%10;
		    n=n/10;
		}
		return x;
	}
\end{code}
\end{answer}
\end{exer}

\newpage
\section{More fun with Algorithms}\label{section:funAlgs}

Let's start with an example of how understanding sets and
set operations can help understand how certain data structures
work.

\begin{exa}\label{exa:TreeSet}
In Java, the \verb+TreeSet+ class is one implementation of a {\em set} that has several methods
with perhaps unfamiliar names, but they do what should be familiar things.
Let's discuss a few of them.\footnote{The method signatures and documentation have been modified from the 
official definition so we can focus on the point at hand.}
Let \verb+A+ and \verb+B+ be \verb+TreeSet+s.
\begin{lenum}
\item  The method \verb+retainAll(TreeSet other)+
``{\em retains only the elements in this TreeSet that are contained in the \verb+other+ TreeSet. 
In other words, removes from this TreeSet all of its elements that are not contained in \verb+other+.}''
It is not too difficult to see that \verb+A.retainAll(B)+ is computing $A\cap B$.\footnote{Technically 
it is doing more than that.  It is storing the result in $A$.  So it is like it is doing $A=A\cap B$,
where $=$ here means assignment, not equals.}
\item  The method \verb+boolean containsAll(TreeSet other)+
``{\em returns true if this set contains all of the elements of \verb+other+ (and false otherwise).}''
Thus, \verb+A.containsAll(B)+ returns true iff $B\subseteq A$.
\item
Even without documentation, it seems likely that \verb+A.size()+ is determining $|A|$.
\item
It also seems likely that \verb+A.isEmpty()+ is determining if $A=\emptyset$.
\end{lenum}
\end{exa}

Here is a more advanced use of the \verb+while+ loop
to compute exponents ($x^n$) more quickly than 
using a simple for loop as was done in Example \ref{exa:powers1}. 

\begin{exa}[Faster Exponentiation] \label{exa:powers2} 
Write an algorithm that is more efficient than the one from 
Example \ref{exa:powers1} to compute $x^n$, 
where $x$ is a given real number and $n$ is a given 
positive integer.
\begin{sol}
The solutions from Example \ref{exa:powers1} requires about $n-1$ multiplications.
Here we give a solution that requires no more than $2\log_2 n$ multiplications, which is significantly better!
We'll start with an example and then describe the process in more general terms.

Let $n=11$.  We can write $n=8+2+1$, which is just expressing $n$ as a sum of powers of $2$.
In other words, we express $n$ in terms of its binary representation.  
Then $x^{11}=x^{8+2+1}=x^8x^2x^1$.  So if we can compute $x^8$, $x^2$, and $x^1$, then we
only need 2 more multiplications to get $x^{11}$.  Notice that we can successively square $x$,
giving the sequence $x \rightarrow x^2 \rightarrow x^4 \rightarrow x^8$ using just $3$ more
multiplications.  So we can compute $x^{11}$ using only $5$ multiplications instead of $10$. 
Notice that $5<2\log_2 11$.

In general, we start by writing $n$ in binary.  We then express $x^n$  in terms of the binary representation
of $n$.  We then successively square $x$ getting a sequence
$$x \rightarrow x^2 \rightarrow x^4 \rightarrow x^8 \rightarrow \cdots \rightarrow x^{2^k},$$
and we stop when $2^k \leq n < 2^{k+1}$.   
As we go, we multiply our answer by the current power of $x$.
Here is an implementation of this idea.
\begin{code}
double power(double x, int n) {
   double ans = 1;
   double pow = x; // Stores x, x^2, x^4, etc.
   int k = n;
   while( k!=0) {
       if(k%2==0) {
           k=k/2;
           pow=pow*pow;
       } else {
           k=k-1;
           ans=ans*pow;
       }
   }
   return ans;
   }
\end{code}
\end{sol}
\end{exa}

\begin{exer}
Trace through the execution of \verb+power(2,23)+. Give a table
with the values of $k$, $pow$, and $ans$ at every step.
How many steps did it take? Was it more efficient than
the algorithm from Example \ref{exa:powers1}? Did it use
as few operations as we promised?

\noindent
{\LARGE
\begin{tabular}{|p{.5in}|p{1.5in}|p{1.5in}|}
\hline
k&pow&ans\\
\hline
&&\\
\hline
&&\\
\hline
&&\\
\hline
&&\\
\hline
&&\\
\hline
&&\\
\hline
&&\\
\hline
&&\\
\hline
&&\\
\hline
&&\\
\hline
\end{tabular}
}
\begin{answer}
Below is a table of the values the algorithm computes.
The final answer of 8388608 is correct!
It took 9 steps as opposed to the 22 required by the 
algorithm from Example \ref{exa:powers1}, so it was
definitely faster. Notice that 
$2 \log_2(23) \approx 2*4.52356=9.04712>9$, 
it works as advertised.

\begin{tabular}{|l|l|l|}
\hline
k&pow&ans\\
\hline
23&2.0&1.0\\
22&2.0&2.0\\
11&4.0&2.0\\
10&4.0&8.0\\
5&16.0&8.0\\
4&16.0&128.0\\
2&256.0&128.0\\
1&65536.0&128.0\\
0&65536.0&8388608.0\\
\hline
\end{tabular}
\end{answer}
\end{exer}



%%%%%%%%%%%%%%%%%%%%
%% The end of this section depends on Partitions
%% and Equivalence Relations...
%%%%%%%%%%%%%%%%%%%%

Now we will see some examples of how
partitions and equivalence relations are relevant to algorithms.
Feel free to skip the rest of this section if you did not
cover Section~\ref{section:parts}.

When creating test cases for an algorithm, you always want to
ensure that you are covering `all of the cases'.  But what does that mean?  It means you are
thinking about how to {\em partition} all of the possible inputs into several sets, where 
the elements in one set are somehow different from those in another set, and are quite a lot
like the other elements in the set.  Let's see an example.

\begin{exa}\label{exa:factPart}
Consider the following function that returns $n!$ if $n\geq 0$, and returns $-1$ if $n<0$
($n!$ is undefined for negative values of $n$, but we have to return something, so why not a negative 
number?)
\begin{code}
	int factorial(int n) {
	    if(n<0)       { return -1; }
	    else if(n==0) { return 1; }
	    else {
	        int fact = 1;
	        for(int i=1;i<=n;i++) {
	            fact = fact*i;
	        }
	        return fact;
	    }
	}
\end{code}
What values of $n$ should we use to test \verb+factorial+?
\begin{sol}
There seems to be three different types of values based on the structure of the code: $0$,
numbers less than $0$, and numbers greater than $0$.  
This suggests partitioning the possible test cases into these
3 classes. If we test at least one number from each of these,
we know we have covered the code in the \verb+if+,
\verb+else if+, and \verb+else+ clauses. Since boundaries can sometimes cause problems, we should include those for each part
as well as an arbitrary value in each part.
In light of this, we might test $0$, $-1$, $-2$, $-10$, $1$, $2$, and $8$.
Since these cover all of the cases, they should provide pretty good evidence of whether
or not \verb+factorial+ is implemented properly.\footnote{But remember, testing never {\em proves}
that code is correct!}
\end{sol}
\end{exa}

It might be helpful to see another example of where equivalence relations and partitions are useful in computer science.

\begin{exa}\label{exa:chance}
Consider a method \verb+setRewardChance(double chance)+ that sets the percentage chance that a player will
be rewarded for completing some task. 
When the chance is set, we may want to take some action based on the value.
For instance, the chance should certainly be non-negative, so we might want to throw some sort of error  
if it is negative. Similarly, the chance should not be above 100 (We are assuming the values are being interpreted 
as whole percentages, so 100 means 100\%). 
We may also want to treat the case of a 0\% chance in a special way. 
Further, chance below 20\%, between 20 and 50\% and above 50\% might all be treated in different ways.
This might lead to code such as the following.
\begin{code}
    setRewardChance(double chance) {
        if(chance<0)         { /* Throw an error */ }
        else if(chance==0)   { /* deal chance==0 */ }
        else if(chance<20)   { /* deal with 0<chance<20 */ }
        else if(chance<50)   { /* deal with 20<= chance<50 */ }
        else if(chance<=100) { /* deal with 50<=chance<=100) */ }
        else                 { /* Throw an error */ }
    }
\end{code}
Notice that a given value of \verb+chance+ will lead us down one of 6 different execution paths.
If we want to write tests for this code, we need to make sure that every path of execution is tested.

So what does this have to do with partitions and equivalence relations? It is simple: The code described
above partitions the set of real numbers into 6 subsets, based on which section of code will be executed:
$$\BBR=(-\infty,0)\cup\{0\}\cup(0,20)\cup[20,50)\cup[50,100]\cup(100,\infty).$$
Therefore, to ensure the tests cover all of the code, we need to choose at least one value from each subset in
the partition. In the words of equivalence classes, we need to choose one or more representatives from
each equivalence class.

Hopefully it is pretty straightforward to see where partitions come into play here. 
But perhaps the concepts related to equivalence relations are more difficult to see.
In this case our equivalence relation is a little more difficult to describe succinctly.  
The most straightforward way is a bit vague: two numbers are related if they are in the same subset.
Although vague, it is somewhat precise (assuming I know that it is referring to the subsets given above). 
Are 40 and 2 related? No, because $40\in[20,50)$, but $2\in(0,20)$.
Are 75 and 98 related? Yes, because $75\in[50,100]$ and $98\in [50,100]$.
\end{exa}

The previous two examples also demonstrates another 
important concept: Although equivalence relations and partitions
are two ways of thinking about the same concept, 
sometimes it is easier to think in terms of one versus the other. 
In Example~\ref{exa:factPart}, it is easiest to think about a
partition of the integers as 
$\BBR=(-\infty,0)\cup\{0\}\cup(0,\infty)$ rather than
formally defining an equivalence relation
on the integers that produces this partition.
You can define the equivalence relation in this case by ``$x$ is related to $y$ if and only if they have the same sign'' if we assume there are
three signs--positive, negative, and no sign (for zero). But that is not the way we usually think about it, so this is at best awkward and at worst incorrect.
Similarly in Example~\ref{exa:chance}, 
the partition $\BBR=(-\infty,0)\cup\{0\}\cup(0,20)\cup[20,50)\cup[50,100]\cup(100,\infty)$ is easy to write down and comprehend,
but it is not so easy to write down a complete and succinct definition of 
the equivalence relation. In fact, I am not even going to try to write
one down in this case---can you think of how to describe this partition
by finishing the phrase ``$x$ is related to $y$ if $\ldots$''? I certainly can't.

On the other hand, 
it is easier to think about congruence modulo $n$ in terms of the 
formal definition of what it means for two integers to be related 
(``$x$ is related to $y$ iff $x\equiv y\pmod n$'').
It is a little more awkward and less informative to write this down as
a partition: $\BBR=[0]\cup[1]\cup[2]\cup\cdots\cup[n-1]$,
especially if you forgot what $[i]$ means.
In any case, it seems to be a little less clear than the definition of the 
equivalence relation given in the first sentence.

In summary, 
in some cases it is easier to talk about what it means for two things to be equivalent, 
and in other cases it is easier to talk about how to partition a set into disjoint subsets. 
But remember that both of these are accomplishing the same thing.\\

In Chapter~\ref{ch:AlgAn}, we will present more sophisticated
algorithms. We wait to present them because it is nice to be able
to evaluate how good an algorithm is. 
We mostly care how fast algorithms are, which is the main topic
of that chapter. But we also sometimes want to know
how much memory an algorithm uses, and how complicated it
is to implement it is also of interest. 
However, those topics are beyond the
scope of this book.

